% !TEX root = ../main.tex
\section{Method}

We developed a framework named \metagent{} for automated literature screening and full-text parameter coding in systematic reviews of infectious disease modeling (Figure~\ref{fig:pipeline}). The proposed \metagent{} comprises three core components: benchmark construction, LLM screening, and full-text parameter coding and extraction. Performance was evaluated separately for the screening and coding components: the screening output was evaluated against the review-derived inclusion lists, while the coding output was evaluated against the estimates reported by the source reviews.

\begin{figure}[!t]
  \centering
  \includegraphics[width=\textwidth]{pipeline/pipeline.pdf}
  \caption{Overview of the \metagent{} pipeline. a) Benchmark construction from published systematic reviews, defining the raw pool of candidate records and the review-derived ground truth. b) LLM-assisted screening based on title-and-abstract metadata, in which a five-criterion screening prompt (disease, population, location, empirical evidence, target parameter) maps each record to Strong, Possible, or Unlikely. c) Full-text parameter coding and extraction comprising data localization (Stage A) and codebook-guided value extraction (Stage B) over the source-review included studies. Evaluation was performed separately for screening and coding outputs and is not shown as a pipeline panel.}
  \label{fig:pipeline}
\end{figure}

\subsection{Benchmark Construction}

This study selected three core epidemiological parameters to evaluate the proposed screening framework: basic reproduction number ($R_0$), serial interval, and fatality, characterizing the transmission potential, transmission speed, and disease severity of an epidemic. We constructed a screening benchmark from published systematic reviews on COVID-19/SARS-CoV-2 and mpox. The benchmark comprises 22 source-review screening tasks, each corresponding to one published systematic review, with a total of 12,758 candidate records. The source reviews used to define the benchmark are listed in Table~\ref{tab:source-reviews}. For each task, we designed a Boolean search query based on the target parameter definitions, using the same search time frame and adhering as closely as possible to the original search strategies reported in the source reviews. PubMed served as the primary data source for retrieving records, which collectively formed the raw pool of candidate articles for screening.

Upon constructing the raw pool, we defined the final inclusion list of each systematic review as the corresponding ground truth (GT) for that task. In this framework, the raw pool delineates the search space of candidate records, whereas the GT identifies the subset of articles that should ultimately be included. Across the benchmark, the aggregated GT consists of 668 articles. It is important to emphasize that the GT in this study does not constitute a rigorous gold standard in the strictest sense. The original search strategies often relied on multiple bibliographic databases and could not be fully reconstructed from the published manuscripts alone, meaning that some GT articles may not appear in our constructed raw pool. The original reviews may themselves contain retrieval omissions or include eligibility constraints that are not fully encoded in the current tasks.

For the coding evaluation, each source review additionally contributed one manually abstracted reference estimate from the source-review report. When available, we selected the reported overall or primary pooled estimate for the target parameter, together with its 95\% confidence interval or reported uncertainty interval. When a source review reported multiple subgroup-, setting-, variant-, or period-specific estimates and did not provide a single overall value, the selected target followed the review's stated primary synthesis question and the included-study evidence available for coding, rather than forcing all reviews into a common estimand. This procedure fixed a transparent review-derived benchmark target for each source review; it was not intended to independently adjudicate whether the source review's own estimate was error-free. The selected reference estimates and uncertainty intervals used in each comparison are listed in Appendix D, enabling readers to audit the target used for each source-review comparison.

\subsection{LLM Screening}

\textbf{Task definition.}
We formulate the task as high-recall first-pass screening, rather than a definitive inclusion/exclusion decision. The objective of the system is to maximally retain GT-positive articles while substantially reducing the number of records requiring manual review, thereby focusing human effort on studies with a higher likelihood of relevance. Final inclusion judgments remain the responsibility of domain experts operating at the full-text level.

\textbf{Screening input.}
Each candidate record was screened on the basis of its bibliographic metadata, comprising the title and the abstract retrieved from PubMed (Figure~\ref{fig:pipeline}b). The same screening prompt and the same Strong, Possible, or Unlikely triage rule were applied uniformly to all records, without auxiliary retrieval or full-text access. When an abstract was unavailable, as occurs for certain preprint records, letters, or older publications with incomplete metadata, the input was reduced to the title alone, while the eligibility criteria remained unchanged; records that could not be confidently excluded from the title alone were retained for downstream review.
% This unified input formulation reflects the operating conditions of a typical database-driven first-pass screening task and isolates the reported screening performance from variability introduced by full-text retrieval or document parsing.

\textbf{Screening prompt.}
Given the title-and-abstract input defined above, the screening prompt structures the assessment around five eligibility criteria:

\begin{itemize}
  \item Disease relevance: the study must explicitly target the infectious disease of interest.
  \item Population relevance: the study should concern human subjects or populations. Purely animal or in vitro studies are considered less relevant.
  \item Location relevance: the study should report data from a real-world geographic area affected by the target disease.
  \item Original empirical evidence: the study should report original empirical data or primary analysis. Reviews, editorials, commentaries, and purely theoretical papers are deemed less relevant.
  \item Parameter relevance: the study should report or estimate the target epidemiological parameter.
\end{itemize}

For each candidate record, the LLM produced a criterion-specific score together with a brief rationale for disease relevance, population relevance, location relevance, original empirical evidence, and target-parameter relevance, and assigned an overall triage label of Strong, Possible, or Unlikely.

The Strong label was assigned to records that provided clear title-and-abstract evidence of the target disease, the target parameter, and original empirical analysis. The Possible label was assigned to records with partial or ambiguous evidence, including sparse records for which relevance could not be confidently excluded from the metadata alone. The Unlikely label was assigned to records that showed no plausible match to the disease, parameter, population, location, or empirical-evidence criteria.

\subsection{Full-Text Parameter Coding and Extraction}

\textbf{Task definition.}
We evaluated full-text parameter coding and extraction as a controlled component-level task (Figure~\ref{fig:pipeline}c). The input to this component was restricted to the studies included in each source review, and the objective was not to make a new eligibility decision but to recover the parameter evidence required for comparison against the review-level epidemiological estimate. This design isolates the coding step from upstream title-and-abstract screening, ensuring that any discrepancy between the coded estimates and the source-review references arises from the coding workflow itself rather than from upstream retrieval errors. Such discrepancies may originate from full-text availability, document parsing, evidence localization, value extraction, or alignment with the synthesis rule of the source review.

\textbf{Full-text processing.}
For each included article, the full text was obtained from open-access sources, with PubMed Central used as the preferred repository. Each retrieved PDF was parsed into a structured text representation that preserved section order, paragraphs, figure captions, and table content where recoverable. This structured representation provided the document input for the subsequent LLM coding stages.

\textbf{Evidence localization and parameter coding.}
Both stages of the coding workflow used Qwen3.6 Plus, the same primary LLM as in the screening analysis, with deterministic decoding and a single inference run per article to support reproducibility. The coding workflow followed the two-stage structure illustrated in the pipeline overview (Figure~\ref{fig:pipeline}c). In Stage A, candidate text passages were localized within the parsed full text, including paragraphs, tables, and figure captions that referred to the target parameter, sample size, denominator, population, study period, geography, uncertainty interval, or estimation method. In Stage B, a parameter-specific codebook was applied to these localized passages to produce structured coding records. Each record captured the parameter type, point estimate, uncertainty interval, sample size or numerator-denominator pair when available, study setting, subgroup definition, and source location within the full text.

Separate codebooks were defined for the three parameter domains: serial interval, reproduction number, and fatality. Serial-interval coding distinguished the target serial interval from related timing quantities, including generation interval, generation time, and incubation period. Reproduction-number coding distinguished the basic reproduction number ($R_0$), the effective reproduction number, and other transmissibility estimates. Fatality coding distinguished case fatality rate (CFR) from infection fatality rate (IFR), hospitalized fatality, ICU fatality, and denominator- or clade-specific fatality estimates. These distinctions were necessary because the same article may report multiple quantities that are not interchangeable for meta-analytic synthesis. The per-article coding records were aggregated to a single review-level estimate for each source review, with continuous parameters pooled by inverse-variance random-effects synthesis.

\subsection{Model Implementation and Comparators}

\textbf{Primary screening model.}
The primary screening model used in the main analysis was Qwen3.6 Plus. This model was selected from a set of candidate screening models on the basis of its overall performance against the predefined operating objective of preserving high recall while reducing the manual screening workload. All candidate models were evaluated on the same source-review benchmark and under the same binary evaluation rule, in which Strong and Possible judgments were treated as records advanced to manual review and Unlikely judgments were treated as records excluded at first pass. Aggregate results for the candidate screening models are reported in Appendix B.

\textbf{Screening baselines.}
Two screening baselines were used. The first was LEADS, a medical literature-mining model for search, screening, and data extraction in evidence synthesis.\cite{ref7} The second was a non-LLM keyword rule based on disease and parameter terms, which provided a simple retrieval-style baseline for first-pass screening.

\textbf{Coding and extraction comparator.}
For the full-text parameter coding and extraction analysis, LEADS was retained only as a coding comparator and was evaluated against the same source-review reference estimates. In contrast to the screening analysis, the coding comparison was not designed as a broad multi-model benchmark; it specifically tested whether the proposed codebook-guided coding and extraction workflow produced review-level estimates that were more closely aligned with the source-review references than those obtained from the LEADS comparator.

\subsection{Evaluation}

\subsubsection{Screening}

The screening performance of the model was evaluated by comparing its outputs against the manually annotated GT. For evaluation purposes, both Strong and Possible relevance judgments were consolidated into a single relevant category, while Unlikely was treated as not relevant. Strong and Possible were both advanced to human review but served different triage functions. Strong records indicated clear title-and-abstract evidence of the target disease, parameter, and original empirical analysis and were prioritized as high-likelihood candidates. Possible records captured incomplete or ambiguous evidence and were retained as a conservative set requiring manual confirmation or full-text assessment. Binary evaluation combined Strong and Possible because both categories remained in the human screening pool. Based on this binary classification, the following four outcomes were defined: true positives (TP), articles classified as relevant that were included in the ground truth; false positives (FP), articles classified as relevant but not included in the ground truth; false negatives (FN), articles classified as not relevant but included in the ground truth; and true negatives (TN), articles classified as not relevant and not included in the ground truth.

From these four categories, we computed recall, precision, workload reduction (WR), and number needed to screen (\nns{}). Recall measures the proportion of GT-relevant articles within the raw pool that were correctly identified by the system. Precision describes the proportion of articles advanced to manual review that are genuinely relevant, thereby quantifying the false-positive burden imposed by a high-recall strategy. WR denotes the proportion of articles classified as not relevant by the system and consequently exempted from manual review. \nns{} represents the average number of articles that must be manually reviewed per eligible study identified; a lower \nns{} indicates greater screening efficiency. Formal metric definitions are provided in Appendix A. Given that the proposed framework is designed as a high-recall screening tool for the title-and-abstract phase, recall and \nns{} were regarded as the primary decision-informative metrics, with precision, F1, and WR reported as secondary operating characteristics.

\subsubsection{Coding}

Coding performance was evaluated as the agreement between LLM-coded estimates and the corresponding source-review reference estimates. For each source review, we computed the absolute difference between the coded point estimate and the reference estimate in the native unit of the review's target parameter (days for serial interval, unitless values for the reproduction number, and percentage points for fatality). Because these parameters are expressed on different scales, MAE was interpreted primarily within parameter domains and reported as parameter-stratified MAE; disease-stratified MAE and the mixed-unit overall MAE were reported only as supplementary descriptive summaries. As a supplementary comparator, the same evaluation procedure was applied to estimates produced by LEADS, and for each source review we determined whether the LLM-coded estimate was closer to the source-review reference than the LEADS estimate. This evaluation should be interpreted as a component-level assessment of coding validity using the source-review included studies as inputs, rather than as an end-to-end measure of performance from search retrieval to final meta-analytic synthesis.
